\section{Conclusion}
\label{sec:conclusion}

Multi-agent pipelines that filter LLM outputs have made the content they inspect safer
while leaving themselves unexamined. We argued that such a pipeline is a multi-agent system
whose own security matters, and we studied two threats with no single-model analogue:
compromised or colluding internal defenders, and second-order prompt injection carried by
the content the judges read. \Aegis answers both by replacing the single trusted
coordinator with Byzantine-robust aggregation of judge verdicts, isolating the untrusted
payload from the judges' instruction channel, hardening each judge, and generalising the
verdict target from jailbreak to task-integrity. We proved, under stated
honest-concentration and margin assumptions, that robust aggregation prevents a Byzantine
minority ($2f{+}2<n$ for selection, $f<n/2$ for median rules) from flipping the aggregate
decision in either direction; we characterised how the guarantee degrades once honest
failures correlate, turning backbone diversity from a tuning choice into a requirement; and
we quantified the linear token/latency cost the committee incurs. We were explicit
throughout that a lower attack success rate under compromise is evidence of graceful
degradation, not a proof of security, and that the empirical figures here are placeholders
fixing a falsifiable protocol.

Two directions follow directly. First, executing the adaptive evaluation of
\cref{sec:experiments}—especially the homogeneous-vs-diverse correlated-failure test and the
security-versus-cost frontier—will tell us whether the preconditions of \cref{thm:integrity}
hold for real hardened LLM judges, which is the crux of whether the approach is practically
meaningful at small $n$. Second, tightening honest concentration without reopening the
inter-agent communication surface—for instance through structured, non-adversarial
consensus among judges—could relax the margin condition and extend tolerance, and is a
question the star-topology architecture is positioned to study. We hope the framing of the
defense pipeline as a first-class attack surface, and the insistence on stating exactly
where robustness stops, are useful beyond the specific architecture proposed here.
